Este capítulo interpreta os resultados à luz da pergunta de pesquisa, separa o que foi demonstrado do que é argumentado e relaciona a arquitetura com o AIXM, com o fluxo atual da informação aeronáutica e com a analogia financeira.

\section{Resposta à pergunta de pesquisa}\label{sec:disc-resposta}

Este trabalho pergunta em que medida uma arquitetura baseada em \textit{event sourcing} torna auditáveis os dados geolocalizados de infraestrutura aeroportuária, e quais são os limites dessa auditabilidade. No estudo de caso do SBSJ, os resultados sustentam a hipótese, e cada uma das cinco propriedades enunciadas no Capítulo~\ref{chap:introducao} encontra respaldo em ao menos um experimento aprovado.
\begin{itemize}
    \item A ausência de rastreabilidade nas publicações atuais é caracterizada no E1, e a rastreabilidade até a causa se confirma no E3, em que o envelope da troca de referencial identifica autor, motivo e causa.
    \item A derivação a partir das premissas é mostrada pelo E2, com os valores publicados recalculados dentro da precisão da AIP.
    \item A reconstrução no tempo é o objeto do E3.
    \item A consistência entre produtos decorre do E5 e se estende aos consumidores no E6. Consistência, aqui, não é correção: produtos que aplicam a mesma regra concordam entre si mesmo se a regra estiver errada.
    \item A evidência de adulteração resulta do E4 e do E6, e se mantém no banco extraído.
\end{itemize}

A resposta, contudo, não é irrestrita. A auditabilidade obtida é documental. A arquitetura garante que se saiba o que foi afirmado, por quem, quando, com base em quais premissas, e que esse registro não foi alterado. Ela não garante que o afirmado seja verdadeiro. Uma medição errada, registrada com autor e causa, permanece errada, embora deixe de ser silenciosa. Essa distinção corresponde à delimitação feita na Introdução, segundo a qual o problema tratado é de rastreabilidade, e não de exatidão.

Os limites, reunidos na Tabela~\ref{tab:disc-demonstrado}, são de dois tipos. Os mais fundamentais dependem de um elemento externo ao registro: a reescrita completa da cadeia só é detectada com uma âncora guardada fora do sistema, e a autoria só é comprovada com assinatura digital. Nenhum mecanismo interno supera essas duas restrições. Os demais decorrem de extensões não implementadas, como a norma de preenchimento e o segundo eixo de tempo, e do escopo da avaliação, restrita a uma pista.

\section{O que é demonstrado e o que é argumentado}\label{sec:disc-demonstrado}

Um dos objetivos específicos deste trabalho é distinguir o que os experimentos demonstram do que é apenas argumentado. A Tabela~\ref{tab:disc-demonstrado} organiza essa distinção por propriedade. Para cada uma, indica o mecanismo que a sustenta, a evidência obtida, a situação do mecanismo e o limite que permanece. A situação tem três níveis: implementado e avaliado nos experimentos, implementado mas não validado no domínio, e apenas argumentado ou proposto.

\begin{table}[htbp]
    \centering
    \caption{Propriedades, mecanismos, evidência e limites.}
    \label{tab:disc-demonstrado}
    \footnotesize
    \renewcommand{\arraystretch}{1.2}
    \begin{tabularx}{\linewidth}{>{\raggedright\arraybackslash}p{2.75cm} >{\raggedright\arraybackslash}X l >{\raggedright\arraybackslash}p{2.2cm} >{\raggedright\arraybackslash}X}
        \toprule
        Propriedade & Mecanismo & Evidência & Situação & Limite \\
        \midrule
        Proveniência causal & envelope com autor, motivo e causa & E1, E3 & implementado e avaliado & autoria declarada, não comprovada \\
        Derivação a partir das premissas & derivados na leitura, pela regra do Anexo~14, com recusa por premissa não declarada & E2 & implementado e avaliado & TODA e LDA conferidas contra premissas lidas da própria AD~2.13; sem regra para barreira de contenção; exatidão do levantamento não avaliada \\
        Posição relativa e remedição & referencial endereçado pelo conteúdo e \textit{seeds} & E2 & implementado e avaliado & uma pista, com dois referentes \\
        Reconstrução no tempo & registro só de acréscimos, com tempo de transação & E3 & implementado e avaliado & histórico documental; relógio simulado nos experimentos; tempo de validade apenas proposto \\
        Evidência de adulteração & cadeia de \textit{hash} sobre forma canônica & E4, E6 & implementado e avaliado & reescrita completa e alteração no fim só com âncora externa \\
        Verificação por terceiro & especificação pública e verificador independente & E5, banco extraído & implementado e avaliado & o arquivo evidencia alterações, mas não as impede; banco extraído publicado gerado com a versão anterior \\
        Consistência entre produtos & produtos como modelos de leitura da mesma fonte, com regra identificada & E5, E6 & implementado e avaliado & consistência não é correção da regra \\
        Leitura separada & tópico, consumidores, limite explícito e grupos & E6 & implementado e avaliado & tópico em memória; sistema de mensagens real, escala e confidencialidade não testados \\
        Geometria derivada & primitivos gerados a partir das \textit{seeds} & -- & implementado, não validado & não comparada às cartas oficiais \\
        Composição do lado ar completo & forma comum do registro & -- & argumentado & não testada com várias pistas; uma cadeia por aeródromo, sem ordem entre aeródromos \\
        Confidencialidade & fluxos separados e controle de acesso & -- & proposto & o tópico geral expõe todos os eventos a quem o lê \\
        Leitura em escala & réplicas, \textit{snapshots} e caches & -- & argumentado & custo não medido \\
        Interoperabilidade & especificação pública e banco extraído & E5, E6 & argumentado & um único publicador \\
        Escrita concorrente & fila de comandos e controle otimista & -- & proposto & bloqueio pessimista, sem teste com várias fontes \\
        Âncora externa & cabeça publicada com a emenda AIRAC & -- & proposto & não publicada nem comparada automaticamente \\
        Norma de preenchimento e bitemporalidade & camadas de campos e tempo de validade & -- & proposto & -- \\
        \bottomrule
    \end{tabularx}
    \par\smallskip{\small Fonte: o autor.}
\end{table}

As afirmações demonstradas valem para o cenário avaliado e foram submetidas a critérios definidos antes da execução. As argumentadas decorrem logicamente da arquitetura, mas não foram postas à prova. As propostas são desenhos que ainda dependem de implementação. A distinção importa porque a arquitetura foi concebida para um aeródromo completo, enquanto a avaliação cobre uma pista. O que se mostra é que a forma do registro sustenta as propriedades nesse recorte. A generalização para o lado ar inteiro depende da hipótese, plausível mas não testada, de que os demais elementos se registram com a mesma forma.

\section{Relação com o AIXM e com a publicação atual}\label{sec:disc-aixm}

A arquitetura não substitui o AIXM. O E1 não emula o modelo AIXM completo, mas apenas a forma como as publicações chegam ao usuário, em edições sucessivas sem registro de causa. O AIXM dispõe de \textit{TimeSlices}, números de sequência e de correção e anotações (Seção~\ref{sec:aixm}). O que falta, e o que a arquitetura acrescenta, está nos três níveis da Seção~\ref{sec:aixm}: no formato, o registro estruturado das premissas e da causa e a evidência de adulteração; na implementação, a retenção obrigatória do histórico; e, no fluxo de publicação, a entrega desse histórico a quem usa os dados.

As duas abordagens são, por isso, complementares. Do ponto de vista do CQRS, uma exportação em AIXM é mais um modelo de leitura do registro, assim como a lente da carta de aeródromo. Cada \textit{TimeSlice} pode ser gerada a partir do estado em um instante, sem que o registro precise adotar o esquema do AIXM. O caminho inverso também é possível. Uma carga inicial a partir de dados em AIXM entra como um grupo de eventos \texttt{BASELINE}, com o provedor como autor. Desse ponto em diante, cada alteração passa a ter causa.

A mesma lógica se aplica ao ciclo AIRAC. A arquitetura não altera a cadência de publicação, mas oferece a ela um complemento natural. A cabeça da cadeia, publicada com cada emenda, funcionaria como âncora externa (Seção~\ref{sec:verificacao}). Quem conservasse as emendas poderia verificar, a qualquer tempo, que o histórico não foi reescrito. Esse ponto toca uma lacuna das normas atuais. O PANS-AIM estabelece requisitos de integridade \cite{icao2018pansaim}, e a DO-200B exige que o usuário tenha um meio de conferir a integridade do que recebe \cite{rtca2015do200b}, mas esse meio protege contra a corrupção na transmissão e no armazenamento. Ele não detecta a reescrita do histórico nem informa a causa de uma mudança.

\section{Impactos operacionais}\label{sec:disc-impactos}

O fluxo atual da informação aeronáutica começa no levantamento e passa pelo operador do aeródromo e pelo serviço de informação aeronáutica até a publicação AIRAC e os usuários. A arquitetura não acrescenta etapas a esse fluxo. O que muda é a forma do dado entre elas.

Cada etapa passa a escrever afirmações, e não valores finais. A equipe de levantamento registra as posições medidas como evento próprio. O operador do aeródromo registra as decisões de projeto, como a extensão de uma zona de parada ou o deslocamento de uma cabeceira, com causa no levantamento ou na obra. O órgão de informação aeronáutica valida e publica, e a publicação passa a ser um corte do registro, e não uma nova cópia dos valores.

Para quem recebe os dados, como operadores aéreos, fornecedores de cartas e a ANAC, três mudanças práticas decorrem disso. A primeira é que a diferença entre duas edições deixa de ser uma comparação manual entre documentos e passa a ser uma consulta, com a causa de cada campo alterado. A segunda é que o recebedor pode verificar a integridade do que recebeu sem depender de quem publicou. A terceira é a redução do trabalho de conciliação entre a AIP e as cartas, pois produtos gerados do mesmo registro para o mesmo instante não divergem por construção.

O próprio caso do SBSJ ilustra o alcance dessas mudanças. Com o registro proposto, a diferença entre a AOC de 2022 e a ADC de 2026 viria acompanhada do evento que a produziu, fosse ele um levantamento, uma obra ou uma correção. A observação da seção AD~2.12 sobre a cabeceira~34 deslocada, discutida na Seção~\ref{sec:req-minimo}, também deixaria de ser ambígua. O deslocamento é uma premissa da direção de pista, e a LDA é calculada a partir dela. No cenário avaliado, a direção~34 afirma deslocamento nulo, com o DECEA como autor e, no motivo, o conflito entre a observação da AD~2.12 e a tabela da AD~2.13. Uma LDA igual à TORA passa, então, a indicar um deslocamento nulo afirmado por alguém, e não uma dúvida silenciosa entre o texto e a tabela. Se o deslocamento não tivesse sido declarado, a LDA seria recusada.

A adoção, por fim, não exige padronização prévia. A arquitetura registra o que existe, com corpo livre, e a norma de preenchimento vem depois, do lado da leitura (Seção~\ref{sec:trabalhos-futuros}). Essa ordem reduz o custo de entrada para órgãos e operadores com dados heterogêneos.

\section{Alcance da analogia financeira}\label{sec:disc-analogia}

A analogia com o livro-razão orienta a arquitetura desde a Introdução, e os resultados permitem delimitar até onde ela vale. Ela vale para o registro e para a auditoria. Nada se apaga, toda correção é um novo lançamento, o valor corrente é derivado do histórico, e dois relatórios do mesmo livro precisam conciliar. O E3 e o E5 confirmam essas propriedades no domínio aeronáutico.

A analogia deixa de valer em dois pontos. O primeiro é a conservação. Na contabilidade, as partidas de um lançamento somam zero, e um lançamento incompleto se denuncia na conciliação. Na informação aeronáutica, os valores derivados resultam de geometria e de regras normativas, sem lei de conservação. Uma pista sem zona de parada é um estado plausível, e uma ASDA igual à TORA não acusa erro algum por si só. As declarações das direções recuperam parte da conferência, pois uma zona de parada declarada e ausente leva à recusa, mas só para as premissas das distâncias declaradas. A atomicidade, que na contabilidade é verificável depois, precisa aqui ser imposta na escrita e respeitada na leitura, o que justifica o mecanismo de grupos.

O segundo ponto é a natureza do objeto registrado. O lançamento contábil constitui o próprio fato, ao passo que o registro aeronáutico descreve um objeto físico que existe independentemente dele (Seção~\ref{sec:bitemporal}). Por isso, um único eixo de tempo basta à contabilidade, mas não esgota as perguntas que se fazem sobre uma pista. A extensão bitemporal decorre, portanto, da diferença entre os domínios, e não de uma deficiência da arquitetura.

\section{Ameaças à validade}\label{sec:disc-validade}

A validade interna é ameaçada, sobretudo, pelo risco de testar o programa contra ele mesmo. A avaliação mitiga esse risco com as três comparações externas descritas na Seção~\ref{sec:validacao}. Uma delas, o verificador independente, encontrou um erro real na forma canônica da biblioteca, o que indica que a comparação externa tem poder de detecção. O risco de ajustar os critérios aos resultados é tratado de duas formas. A principal é a derivação das margens a partir da precisão de publicação e da exatidão do gabarito, que não dependem dos valores obtidos. A segunda é o registro das verificações no histórico do repositório antes da execução registrada, sem alteração posterior. Esse registro mostra a ordem dos \textit{commits}, mas não exclui execuções locais anteriores, e por isso não equivale a um pré-registro formal.

A validade externa é a mais limitada. A avaliação usa um único estudo de caso, com uma pista, e parte do cenário é simulada. As premissas e os valores de comparação são reais, mas a remedição e as restrições não ocorreram. Os resultados mostram que a arquitetura é capaz de sustentar as propriedades, e não que as sustentará em qualquer aeródromo ou volume.

Quanto à validade de construto, a auditabilidade foi operacionalizada como cinco propriedades verificáveis. Essa operacionalização não esgota o conceito. A autoria é registrada, mas não autenticada, e uma auditoria regulatória exigiria também a comprovação de identidade. A medição da qualidade dos dados, como a exatidão de um levantamento, está fora do construto por decisão de escopo.

\section{Natureza da contribuição}\label{sec:disc-contribuicao}

Na classificação de \citeonline{gregor2013}, a contribuição deste trabalho é uma exaptação, isto é, a aplicação de soluções conhecidas em um domínio a um problema de outro domínio. Nenhum dos mecanismos empregados é novo isoladamente. O \textit{event sourcing}, o encadeamento por \textit{hash}, o referenciamento linear e o CQRS têm literatura e uso consolidados. A contribuição está na sua composição para os dados de infraestrutura aeroportuária e na adaptação que esse domínio exige. Exemplos dessa adaptação são as posições relativas a referentes, os derivados calculados por regras normativas, a atomicidade imposta na escrita e a distinção entre o registro e o objeto físico. A Tabela~\ref{tab:origem} indica, elemento a elemento, o que foi adotado, adaptado ou proposto.

Os mecanismos organizam-se em torno de três propriedades. A primeira é a proveniência causal, isto é, a capacidade de dizer por que um valor mudou, quem o afirmou e de qual evento a mudança decorre. A segunda é a derivação a partir de premissas retidas, pela qual os valores publicados são recalculados, e não guardados, e uma nova medição se propaga sem reescrita. A terceira é a auditabilidade do histórico, que permite reconstruir qualquer instante e verificar, sem confiar em quem publica, que o que se recebeu não foi alterado. Os demais elementos, como o tópico, os consumidores, os caches e o banco extraído, são os meios pelos quais essas propriedades chegam a quem lê.

A contribuição materializa-se em quatro artefatos. O primeiro é a arquitetura, descrita de forma independente da implementação. O segundo é a biblioteca de código aberto, que mostra que a arquitetura é implementável com tecnologia comum. O terceiro é o banco extraído do SBSJ, publicado como conjunto de dados aberto, que permite a terceiros reproduzir parte da avaliação sem a biblioteca. O quarto é o método de avaliação por falseamento, com critérios definidos antes da execução, reprodutível a partir dos repositórios. O Capítulo~\ref{chap:conclusao} retoma esses resultados e indica os trabalhos futuros.